\section{Validation Mode}
\label{sec:ValidationMode}

\subsection{Purpose}

Validation mode runs a not-yet-released implementation of a process task against
real input, alongside the current (or prior) version, and reports where the new
version produces the same output and where it diverges. It answers a concrete
release question: \emph{if I deploy this candidate task, what changes downstream?}
Sometimes an identical input$\rightarrow$output is expected (a refactor); sometimes
divergence is intended (a functional correction). Validation mode does not judge ---
it produces the evidence and lets a human decide.

\subsection{Per-task isolation}

Validation is handled \emph{per task}. Each candidate task is fed the production
input \emph{for that task} --- the prior version's output of its predecessor --- not
the candidate's own reworked output. Candidate outputs are therefore never chained
together. A divergence introduced by a modified early task cannot cascade into and
contaminate the comparison of later tasks: every task is compared in isolation
against production for the identical shared input.

\subsection{Shadow worker and the dedicated index}

For each plugin task, scaffolding with \texttt{--validation} generates an additional
\emph{shadow worker} beside the production worker. The shadow worker:
\begin{itemize}
\item consumes the same production input topic through a \emph{dedicated consumer
  group} (\texttt{\%processId\%\_\\\%taskId\%\_validation}), so the production consumer
  offset --- the current input index --- is never advanced or disturbed;
\item starts at \texttt{latest} by default (live-concurrent shadowing) and can be
  pointed at recent history via \texttt{DURGA\_VALIDATION\_OFFSET\_RESET} (a bounded
  replay near the tail --- not from the very beginning, since old input may not be
  representative);
\item runs the candidate implementation with side effects suppressed;
\item writes nothing to the task output topic, \texttt{process-events}, or the
  dead-letter queue. Its output is diverted to the shared
  \texttt{validation-candidate-outputs} topic as a \\ \texttt{ValidationCandidateOutput}
  record carrying the shared input, the candidate output, disposition, idempotency
  key, and any error.
\end{itemize}

Live-concurrent and historic replay differ only in where the dedicated consumer
starts; the matching mechanism is identical in both.

\subsection{Side-effect suppression}

Because the shadow worker never publishes to the task output topic,
\texttt{process-events}, or the DLQ, downstream tasks never observe a candidate
result. For handle-aware (\texttt{materialize}) execution, \\
\texttt{PluginExecutionSupport.executeSandboxed} runs the transform but writes no
object to the store: it returns a synthetic handle whose content hash still reflects
the real output, so the meaningful part stays comparable while volatile fields (uri,
timestamps) are excluded by comparison ignore paths. Version~1 targets pure transform
(plugin) tasks; tasks that perform their own external effects (send tasks, external
calls, object-store writes inside a plugin body) are out of scope for validation and
must not be shadowed.

\subsection{Comparison and report}

The monitor's \texttt{ValidationTopology} pairs each candidate output against the
prior/production output for the same input --- the \texttt{ACTIVITY\_COMPLETED}
lifecycle event --- keyed identically on
\texttt{processId:\\activityId:processInstanceId}. The two sides are merged into a
single aggregate so the comparison is robust to arrival order: in live mode the
candidate may arrive before or after the production output. A candidate seen before
any prior output yields \texttt{PRIOR\_MISSING}, superseded when the prior output
later arrives.

Comparison is a structural, order-sensitive JSON diff (\texttt{JsonComparison}) with
configurable \emph{ignore paths} (dot-notation, with \texttt{*} and \texttt{[*]}
wildcards; set via \texttt{durga.validation.ignore.paths}) so volatile fields such as
generated timestamps or identifiers do not register as divergence. Each comparison
produces a \texttt{ValidationResult} with one of four outcomes:

\begin{description}
\item[\texttt{EQUAL}] candidate output identical to prior after ignore paths.
\item[\texttt{DIFF}] candidate output differs; the located differences (path, kind,
  prior value, candidate value) are attached.
\item[\texttt{PRIOR\_MISSING}] no production output found for the shared input.
\item[\texttt{CANDIDATE\_ERROR}] the candidate failed or requested a non-success
  error strategy.
\end{description}

Results are written to \texttt{validation-results} and materialized in the
\texttt{validation-results-store} global table for query.

\subsection{Topics and naming}

\begin{itemize}
\item \texttt{\%processId\%\_\%taskId\%\_validation} --- dedicated consumer group for
  the shadow worker (reads the production input topic).
\item \texttt{validation-candidate-outputs} --- shared topic carrying candidate
  outputs from every shadow worker (Java and Rust).
\item \texttt{validation-results} --- comparison results, materialized in
  \texttt{validation-results-store}.
\end{itemize}

\subsection{Reporting interfaces}

The monitor exposes the report under \texttt{/api/validation/}:
\begin{itemize}
\item \texttt{/api/validation/summary?processId=\textit{id}} --- per-task counts
  (equal / diff / prior-missing / candidate-error / total); \texttt{processId} is
  optional (omit it for all processes),
\item \texttt{/api/validation/results?processId=\&taskId=\&status=} --- individual
  comparisons, optionally filtered,
\item \texttt{/api/validation/instances/\{instanceId\}} --- all validated
  tasks for one instance.
\end{itemize}

The dashboard adds a \emph{Validation Report} panel to the per-process drill-down: a
per-task summary table and, per instance, an expandable comparison showing the shared
input, prior output, candidate output, and a highlighted diff table. Prometheus
metrics \texttt{durga\_validation\_total}, \texttt{durga\_validation\_diff}, and
\texttt{durga\_validation\_candidate\_error} are exported per
\texttt{process\_id,task\_id} on \\ \texttt{/api/metrics}.

\subsection{Java and Rust parity}

Both code-generation targets emit shadow workers. The
\texttt{ValidationCandidateOutput} wire record is byte-compatible across targets
(camelCase fields matching the Java record), so a mixed Java/Rust fleet feeds the same
\texttt{validation-candidate-outputs} topic and the same monitor comparator. The Rust
runtime crate (\filepath{durga-rust}) provides the matching
\texttt{ValidationCandidateOutput} type and a \texttt{plan\_validation\_output}
function; generated Rust projects add a \texttt{\%taskId\%\_validation} binary per
plugin task using \texttt{run\_validation\_worker} with the same dedicated-group and
offset-reset behaviour.

\subsection{Comparison model summary}

\begin{longtable}{|>{\raggedright\arraybackslash}p{0.35\textwidth}|>{\raggedright\arraybackslash}p{0.57\textwidth}|}
\hline
\rowcolor[gray]{0.9}\textbf{Type} & \textbf{Role} \\
\hline
\endfirsthead
\hline
\rowcolor[gray]{0.9}\textbf{Type} & \textbf{Role} \\
\hline
\endhead
\texttt{ValidationCandidateOutput} & Shadow worker $\rightarrow$ comparator wire record: shared input, candidate output, disposition, error. \\
\texttt{JsonComparison} & Structural JSON diff with ignore-path normalization. \\
\texttt{ValidationResult} & Comparator $\rightarrow$ report record: outcome plus located differences and both payloads. \\
\texttt{ValidationSummary} & Per-task aggregate counts surfaced to the API, dashboard, and metrics. \\
\end{longtable}
